% GENERATED FILE. Edit canonical lessons, then run npm run generate:books.
% canonical-source-hash: fnv1a64:7c1dc3d6a28176a9
% canonical-lessons: ZH-C221-xueren, ZH-W221-xue-snow, ZH-C221-huaxue, ZH-C221-bing, ZH-C221-lingxia, ZH-C221-shoutao

\chapter{Snow and Ice}
\label{ch:zh-snow-and-ice}
\hlchaptermodality{221}

\begin{chapteropening}
\textbf{By the end of this chapter:} \emph{I can talk about a snowman, skiing, ice, being below zero and gloves.}

Complete the last lesson of chapter 221: I can talk about a snowman, skiing, ice, being below zero and gloves.
\end{chapteropening}

\section[\texorpdfstring{\zh{雪人} (xuěrén)}{雪人 (xuěrén)}]{\zh{雪人} (xuěrén) — snowman}

\label{lesson:ZH-C221-xueren}



\begin{quote}
Before the new one: say the Mandarin for to skip with a rope, then the Mandarin for to do a fitness routine.
\end{quote}

\subsection*{You'll want to know: \zh{雪人}}
\textbf{\zh{雪人}} — \emph{xuěrén} — \textquotedblleft{}snowman\textquotedblright{}.

\begin{quote}
\textbf{\zh{门口有一个雪人}} \emph{ménkǒu yǒu yí ge xuěrén} — There's a snowman by the door.
\end{quote}

The new character, in pieces: \zh{雪} — eleven strokes in two pieces: \zh{雨} pressed flat on top, eight strokes, and under it three bars joined on the right. The next lesson puts it in your hand.

\begin{cousinweb}
\textbf{\zh{雪人}} (\emph{xuěrén}) is built from \zh{雪} 'snow' and \zh{人} 'person', both inherited from Middle Chinese, \emph{sjwet} and \emph{nyin} (Guangyun, in Baxter's transcription), which glosses the first 'frozen rain'.
\end{cousinweb}

\begin{grammarlens}[title={What you've built}]
One more word for describing things.
\end{grammarlens}

\subsection*{Your turn}
\begin{itemize}
\raggedright
  \item \emph{Say it:} \emph{xuěrén}
  \item \emph{Say it:} \emph{xuěrén}, once more
  \item \emph{Say it:} \emph{ménkǒu yǒu yí ge xuěrén} — There's a snowman by the door.
  \item \emph{Recall:} say the Mandarin for to skip with a rope, then the Mandarin for to do a fitness routine, then say \emph{xuěrén} again
\end{itemize}

\begin{tcolorbox}[breakable,colback=teal!4,colframe=teal!35!black,title={Before you move on}]
What does \zh{雪人} mean? (\textquotedblleft{}Snowman\textquotedblright{}.) Say it once more.
\end{tcolorbox}
\section[\texorpdfstring{\zh{雪} (xuě)}{雪 (xuě)}]{\zh{雪} — \zh{雨} on top, three bars below}

\label{lesson:ZH-W221-xue-snow}



\begin{quote}
The word first, then the character inside it.

\begin{itemize}
\raggedright
  \item \emph{Recall:} say \emph{xuěrén} — \textbf{R1}, one lesson back, before it decays
\end{itemize}

\begin{quote}
\zh{雨} \zh{白}
\end{quote}

\emph{Read it:} each one aloud, with nothing said first — \emph{yǔ}, \emph{bái}
\end{quote}

\subsection*{Script — \zh{雨} on top, three bars below}
\hlblockfigure{\detokenize{figures/ZH-W221-xue-snow-filmstrip.pdf}}{How \zh{雪} is written, stroke by stroke}

\begin{quote}
\zh{雪}
\end{quote}

\textbf{Two pieces, top before bottom.} Count the pieces, not the strokes.

\textbf{On top, \zh{雨}, eight strokes,} pressed flat: a short bar; a short stroke down on the left; a bar that runs right and turns down with a hook; a vertical through the middle; then four dots, two on the left and two on the right.

\textbf{Under it, three strokes.} A bar that runs right and turns down, a bar across the middle, and a bar along the foot, all three meeting the stroke on the right.

8 + 3 = eleven.

\subsection*{Your turn}
\begin{itemize}
\raggedright
  \item \emph{Write it:} \zh{雪}, all eleven, with the model in front of you
  \item \emph{Cover:} write it again from memory, and count the pieces, not the strokes
\end{itemize}

\begin{tcolorbox}[breakable,colback=teal!4,colframe=teal!35!black,title={Before you move on}]
How many strokes in \zh{雪}? (\textbf{Eleven}.) What is on top? (\textbf{\zh{雨}, pressed flat}.) How many bars at the foot? (\textbf{Three}.)

Source: \href{https://raw.githubusercontent.com/chanind/hanzi-writer-data/68d10a4b21150cae5e1ebbd223eed289cf32d90c/data/\%E9\%9B\%AA.json}{Hanzi Writer Data \zh{雪}, pinned stroke paths and medians}.
\end{tcolorbox}
\section[huáxuě]{huáxuě — to ski}

\label{lesson:ZH-C221-huaxue}



\begin{quote}
Before the new one: say the Mandarin for to do a fitness routine, then the Mandarin for a snowman.
\end{quote}

\subsection*{You'll want to know: huáxuě}
\emph{huáxuě} — \textquotedblleft{}to ski\textquotedblright{}.

\begin{quote}
\emph{nǐ huì huáxuě ma} — Can you ski?
\end{quote}

Its characters come later; for now it is a word for the ear.

\begin{cousinweb}
\emph{huáxuě} is built from \emph{huá} 'slippery' and \zh{雪} 'snow', both inherited from Middle Chinese, \emph{hweat} and \emph{sjwet} (Guangyun, in Baxter's transcription). The dictionary's gloss for the first is 'keen' or 'smooth', not 'slippery', so how it came to that sense is not settled.
\end{cousinweb}

\begin{grammarlens}[title={What you've built}]
One more word for describing things.
\end{grammarlens}

\subsection*{Your turn}
\begin{itemize}
\raggedright
  \item \emph{Say it:} \emph{huáxuě}
  \item \emph{Say it:} \emph{huáxuě}, once more
  \item \emph{Say it:} \emph{nǐ huì huáxuě ma} — Can you ski?
  \item \emph{Recall:} say the Mandarin for to do a fitness routine, then the Mandarin for a snowman, then say \emph{huáxuě} again
\end{itemize}

\begin{tcolorbox}[breakable,colback=teal!4,colframe=teal!35!black,title={Before you move on}]
What does huáxuě mean? (\textquotedblleft{}To ski\textquotedblright{}.) Say it once more.
\end{tcolorbox}
\section[bīng]{bīng — ice}

\label{lesson:ZH-C221-bing}



\begin{quote}
Before the new one: say the Mandarin for a snowman, then the Mandarin for to ski.
\end{quote}

\subsection*{You'll want to know: bīng}
\emph{bīng} — \textquotedblleft{}ice\textquotedblright{}.

\begin{quote}
\emph{bīngxiāng lǐ yǒu bīng} — There's ice in the fridge.
\end{quote}

Its characters come later; for now it is a word for the ear.

\begin{cousinweb}
\emph{bīng} is inherited from Middle Chinese \emph{ping} (Guangyun, in Baxter's transcription), glossed 'water frozen'.
\end{cousinweb}

\begin{grammarlens}[title={What you've built}]
One more word for describing things.
\end{grammarlens}

\subsection*{Your turn}
\begin{itemize}
\raggedright
  \item \emph{Say it:} \emph{bīng}
  \item \emph{Say it:} \emph{bīng}, once more
  \item \emph{Say it:} \emph{bīngxiāng lǐ yǒu bīng} — There's ice in the fridge.
  \item \emph{Recall:} say the Mandarin for a snowman, then the Mandarin for to ski, then say \emph{bīng} again
\end{itemize}

\begin{tcolorbox}[breakable,colback=teal!4,colframe=teal!35!black,title={Before you move on}]
What does bīng mean? (\textquotedblleft{}Ice\textquotedblright{}.) Say it once more.
\end{tcolorbox}
\section[língxià]{língxià — below zero}

\label{lesson:ZH-C221-lingxia}



\begin{quote}
Before the new one: say the Mandarin for to ski, then the Mandarin for ice.
\end{quote}

\subsection*{You'll want to know: língxià}
\emph{língxià} — \textquotedblleft{}below zero\textquotedblright{}.

\begin{quote}
\emph{jīntiān língxià wàibian hěn lěng} — It's below zero today; it's freezing outside.
\end{quote}

Its characters come later; for now it is a word for the ear.

\begin{cousinweb}
\emph{língxià} is built from \emph{líng} 'zero' and \zh{下} 'below', both inherited from Middle Chinese, \emph{leng} and \emph{haeX} (Guangyun, in Baxter's transcription). The dictionary glosses the first 'to fall' and 'a gentle rain', so the sense zero came later, and how is not settled.
\end{cousinweb}

\begin{grammarlens}[title={What you've built}]
One more word for describing things.
\end{grammarlens}

\subsection*{Your turn}
\begin{itemize}
\raggedright
  \item \emph{Say it:} \emph{língxià}
  \item \emph{Say it:} \emph{língxià}, once more
  \item \emph{Say it:} \emph{jīntiān língxià wàibian hěn lěng} — It's below zero today; it's freezing outside.
  \item \emph{Recall:} say the Mandarin for to ski, then the Mandarin for ice, then say \emph{língxià} again
\end{itemize}

\begin{tcolorbox}[breakable,colback=teal!4,colframe=teal!35!black,title={Before you move on}]
What does língxià mean? (\textquotedblleft{}Below zero\textquotedblright{}.) Say it once more.
\end{tcolorbox}
\section[shǒutào]{shǒutào — gloves}

\label{lesson:ZH-C221-shoutao}



\begin{quote}
Before the new one: say the Mandarin for ice, then the Mandarin for below zero.
\end{quote}

\subsection*{You'll want to know: shǒutào}
\emph{shǒutào} — \textquotedblleft{}gloves\textquotedblright{}.

\begin{quote}
\emph{wàibian hěn lěng nǐ dài shǒutào le ma} — It's cold out. Have you got your gloves?
\end{quote}

Its characters come later; for now it is a word for the ear.

\begin{cousinweb}
\emph{shǒutào} is built from \zh{手} 'hand' and \emph{tào} 'a cover', both inherited from Middle Chinese, \emph{syuwX} and \emph{thawX} (Guangyun, in Baxter's transcription). The dictionary glosses the second 'long', so how it came to mean a cover is not settled.
\end{cousinweb}

\begin{grammarlens}[title={What you've built}]
One more word for describing things.
\end{grammarlens}

\subsection*{Your turn}
\begin{itemize}
\raggedright
  \item \emph{Say it:} \emph{shǒutào}
  \item \emph{Say it:} \emph{shǒutào}, once more
  \item \emph{Say it:} \emph{wàibian hěn lěng nǐ dài shǒutào le ma} — It's cold out. Have you got your gloves?
  \item \emph{Recall:} say the Mandarin for ice, then the Mandarin for below zero, then say \emph{shǒutào} again
\end{itemize}

\begin{tcolorbox}[breakable,colback=teal!4,colframe=teal!35!black,title={Before you move on}]
What does shǒutào mean? (\textquotedblleft{}Gloves\textquotedblright{}.) Say it once more.
\end{tcolorbox}
